\documentclass[12pt]{extarticle}

% ===== Encodage & langue =====
\usepackage{fontspec}
\usepackage{polyglossia}
\setdefaultlanguage{french}

% ===== Police Verdana (compiler avec XeLaTeX ou LuaLaTeX) =====
\setsansfont{Verdana}
\renewcommand{\familydefault}{\sfdefault}

% ===== Interligne =====
\usepackage{setspace}
\setstretch{1.1}

% ===== Marges réduites =====
\usepackage{geometry}
\geometry{a4paper, top=1.8cm, bottom=1.8cm, left=2cm, right=2cm}

% ===== Mathématiques =====
\usepackage{amsmath}

% ===== Espacement réduit autour des équations =====
\setlength{\abovedisplayskip}{4pt}
\setlength{\belowdisplayskip}{4pt}
\setlength{\abovedisplayshortskip}{2pt}
\setlength{\belowdisplayshortskip}{2pt}

% ===== Espacement réduit des sections =====
\usepackage{titlesec}
\titlespacing*{\section}{0pt}{1.2ex}{0.6ex}

% ===== Titre du document =====
\title{\textbf{Correction du contrôle séquence 1 version 3}}
\author{}
\date{}

\begin{document}
\maketitle

% --------------------------------------------------
\section*{Exercice 1}
\begin{align*}
15 - 3 - 4 + 25 - 5 &= 12 - 4 + 25 - 5 \\
                     &= 8 + 25 - 5 \\
                     &= 33 - 5 \\
                     &= 28
\end{align*}

% --------------------------------------------------
\section*{Exercice 2}

\begin{align*}
\textbf{a)}
7 + 5 \times 4 &= 7 + 20 \\
               &= 27
\end{align*}

\textbf{b)}
\begin{align*}
46 - 8 \times 4 &= 46 - 32 \\
                &= 14
\end{align*}

\textbf{c)}
\begin{align*}
43 + 18 \div 6 &= 43 + 3 \\
               &= 46
\end{align*}

\textbf{d)}
\begin{align*}
8 + 7 \times 5 + 4 - 8 \div 2 &= 8 + 35 + 4 - 4 \\
                              &= 43
\end{align*}

% --------------------------------------------------
\section*{Exercice 3}
\begin{align*}
3 \times 8 + 3 \times 7 &= 24 + 21 \\
                        &= 45
\end{align*}

La dépense est de 45~€.

% --------------------------------------------------
\section*{Exercice 4}

\textbf{a)} $36 - (7 + 9) \div 2$
\begin{align*}
&= 36 - 16 \div 2 \\
&= 36 - 8 \\
&= 28
\end{align*}

\textbf{b)} $(5 + 9) \times 4 - 8 + 6$
\begin{align*}
&= 14 \times 4 - 8 + 6 \\
&= 56 - 8 + 6 \\
&= 54
\end{align*}

\textbf{c)} $(23 - (6 + 8)) \times 5$
\begin{align*}
&= (23 - 14) \times 5 \\
&= 9 \times 5 \\
&= 45
\end{align*}

\textbf{d)} $6 \times (27 - (8 + 14)) - 5$
\begin{align*}
&= 6 \times (27 - 22) - 5 \\
&= 6 \times 5 - 5 \\
&= 30 - 5 \\
&= 25
\end{align*}

% --------------------------------------------------
\section*{Exercice 5}
\begin{align*}
16 - 6 \times 2 &= 16 - 12 \\
                &= 4
\end{align*}

Une gomme coûte 4~€.

% --------------------------------------------------
\section*{Exercice bonus 1}

\textbf{a)} $9 \times 5 + 12$

\textbf{b)} $6 \times (8 + 3)$

% --------------------------------------------------
\section*{Exercice bonus 2}

\textbf{a)} La somme du produit de 7 par 5 et de 8.

\textbf{b)} Le produit de la somme de 8 et de 3 par 6.

\end{document}